
\documentclass[11pt,letterpaper]{article}

% ============================================================
% COSMIC UNIVERSALISM WEBSITE MASTER IMPLEMENTATION MANUAL
% Compile with: pdfLaTeX
% Current stable website integration verified: July 23, 2026
% ============================================================

% ------------------------------------------------------------
% Core packages
% ------------------------------------------------------------
\usepackage[margin=0.78in,headheight=16pt]{geometry}
\usepackage[T1]{fontenc}
\usepackage[utf8]{inputenc}
\usepackage{lmodern}
\usepackage[scaled=0.94]{helvet}
\usepackage{microtype}
\usepackage{parskip}
\usepackage{setspace}
\usepackage{enumitem}
\usepackage{booktabs}
\usepackage{tabularx}
\usepackage{array}
\usepackage{longtable}
\usepackage{xcolor}
\usepackage{graphicx}
\usepackage{fancyhdr}
\usepackage{lastpage}
\usepackage{titlesec}
\usepackage{hyperref}
\usepackage{xurl}
\usepackage{bookmark}
\usepackage{listings}
\usepackage[most]{tcolorbox}
\usepackage{amssymb}
\usepackage{amsmath}
\usepackage{pifont}
\usepackage{needspace}
\usepackage{caption}
\usepackage{tocloft}

% ------------------------------------------------------------
% Document typography
% ------------------------------------------------------------
\renewcommand{\familydefault}{\sfdefault}
\setstretch{1.08}
\setlength{\parindent}{0pt}
\setlength{\parskip}{0.65em}
\raggedbottom
\urlstyle{same}
\setlength{\cftsecnumwidth}{3.3em}
\setlength{\cftsubsecnumwidth}{4.2em}

% ------------------------------------------------------------
% Cosmic Universalism color palette
% ------------------------------------------------------------
\definecolor{CUNavy}{HTML}{050812}
\definecolor{CUNavyTwo}{HTML}{080D1A}
\definecolor{CUSurface}{HTML}{0C1424}
\definecolor{CUGold}{HTML}{DDBF72}
\definecolor{CUGoldLight}{HTML}{FFF8DD}
\definecolor{CUCyan}{HTML}{35CDEB}
\definecolor{CUCyanLight}{HTML}{78E8FF}
\definecolor{CUBlue}{HTML}{668ED1}
\definecolor{CUBlueLight}{HTML}{9EB9EA}
\definecolor{CUAmber}{HTML}{D99A32}
\definecolor{CURed}{HTML}{B84C4C}
\definecolor{CUGreen}{HTML}{2C8A65}
\definecolor{CULight}{HTML}{F3F6FA}
\definecolor{CUBorder}{HTML}{CCD5E0}
\definecolor{CUText}{HTML}{18212E}
\definecolor{CUMuted}{HTML}{5C6878}
\definecolor{CUCodeBackground}{HTML}{F6F8FB}
\definecolor{CUCodeBorder}{HTML}{BFCADA}

% ------------------------------------------------------------
% Project metadata
% ------------------------------------------------------------
\newcommand{\ProjectTitle}{Cosmic Universalism CU-Time}
\newcommand{\DocumentTitle}{Expanded Outputs and Method Implementation Manual}
\newcommand{\DocumentSubtitle}{Full Cosmic Breath Coordinate, Observable-Universe Reference Outputs, and Controlled Astro Integration}
\newcommand{\RepositoryName}{CosmicUniversalismStatement}
\newcommand{\RepositoryURL}{https://github.com/willmaddock/CosmicUniversalismStatement}
\newcommand{\ProtectedBranch}{main}
\newcommand{\StableBranch}{website}
\newcommand{\ProposedBranch}{feature/cu-time-expanded-outputs}
\newcommand{\WebsiteDirectory}{website/}
\newcommand{\TargetRoute}{/CosmicUniversalismStatement/cu-time/}
\newcommand{\CurrentWebsiteCommit}{e1c82bf}
\newcommand{\CurrentWebsiteFullCommit}{e1c82bff1d37732f52b03fa4734cf7e4120b71e9}
\newcommand{\CurrentIntegrationTag}{website-v1.1-cosmic-breath}
\newcommand{\CurrentIntegrationCommit}{7ac87167c51aad9ae7a31e17d1b2127fc069e15c}
\newcommand{\PriorCuTimeTag}{website-v1.1-cu-time-release}
\newcommand{\PriorCuTimeCommit}{3fa4a37}
\newcommand{\DocumentVersion}{1.1}
\newcommand{\VerificationDate}{July 23, 2026}
\newcommand{\AuthorName}{William Maddock}

% ------------------------------------------------------------
% Hyperlinks and PDF metadata
% ------------------------------------------------------------
\hypersetup{
    colorlinks=true,
    linkcolor=CUBlue,
    urlcolor=CUBlue,
    citecolor=CUBlue,
    pdftitle={Cosmic Universalism CU-Time Expanded Outputs and Method Implementation Manual},
    pdfauthor={William Maddock},
    pdfsubject={Controlled expansion of the Astro CU-Time page with observable-universe reference outputs and CU-governed interpretation},
    pdfkeywords={Cosmic Universalism, CU-Time, Cosmic Breath, sub-ZTOM, ZTOM,
                 Observable Universe, Big Bang reference, Astro, TypeScript,
                 Decimal.js, Gregorian UTC, Julian Day Number, Vitest}
}

% ------------------------------------------------------------
% Header and footer
% ------------------------------------------------------------
\pagestyle{fancy}
\fancyhf{}
\fancyhead[L]{\small\textbf{Cosmic Universalism CU-Time}}
\fancyhead[R]{\small\DocumentTitle}
\fancyfoot[L]{\small Version \DocumentVersion}
\fancyfoot[C]{\small \thepage\ of \pageref{LastPage}}
\fancyfoot[R]{\small Scope: \texttt{Expanded Outputs}}
\renewcommand{\headrulewidth}{0.5pt}
\renewcommand{\footrulewidth}{0.5pt}

% ------------------------------------------------------------
% Section formatting
% ------------------------------------------------------------
\titleformat{\section}
  {\Large\bfseries\color{CUNavy}}
  {\thesection}
  {0.8em}
  {}
  [\vspace{0.2em}\color{CUCyan}\titlerule]

\titleformat{\subsection}
  {\large\bfseries\color{CUNavyTwo}}
  {\thesubsection}
  {0.7em}
  {}

\titleformat{\subsubsection}
  {\normalsize\bfseries\color{CUBlue}}
  {\thesubsubsection}
  {0.6em}
  {}

\titlespacing*{\section}{0pt}{2.1em}{0.9em}
\titlespacing*{\subsection}{0pt}{1.5em}{0.65em}
\titlespacing*{\subsubsection}{0pt}{1.1em}{0.35em}

% ------------------------------------------------------------
% Lists and tables
% ------------------------------------------------------------
\setlist[itemize]{leftmargin=1.6em,itemsep=0.3em,topsep=0.3em}
\setlist[enumerate]{leftmargin=2em,itemsep=0.4em,topsep=0.35em}
\newcommand{\checkbox}{\(\square\)}
\newcommand{\checked}{\ding{51}}
\newcolumntype{Y}{>{\raggedright\arraybackslash}X}
\newcolumntype{L}[1]{>{\raggedright\arraybackslash}p{#1}}

% ------------------------------------------------------------
% General information boxes
% ------------------------------------------------------------
\newtcolorbox{purposebox}{
    enhanced, breakable, colback=CULight, colframe=CUBlue,
    boxrule=0.8pt, arc=2mm, left=3mm, right=3mm, top=2.5mm, bottom=2.5mm,
    title=\textbf{Purpose}, coltitle=white, colbacktitle=CUBlue
}
\newtcolorbox{checkpointbox}{
    enhanced, breakable, colback=CUCyan!6, colframe=CUCyan,
    boxrule=0.8pt, arc=2mm, left=3mm, right=3mm, top=2.5mm, bottom=2.5mm,
    title=\textbf{Checkpoint}, coltitle=CUNavy, colbacktitle=CUCyanLight
}
\newtcolorbox{safetybox}{
    enhanced, breakable, colback=CUAmber!8, colframe=CUAmber,
    boxrule=0.8pt, arc=2mm, left=3mm, right=3mm, top=2.5mm, bottom=2.5mm,
    title=\textbf{Safety Rule}, coltitle=white, colbacktitle=CUAmber
}
\newtcolorbox{warningbox}{
    enhanced, breakable, colback=CURed!6, colframe=CURed,
    boxrule=0.8pt, arc=2mm, left=3mm, right=3mm, top=2.5mm, bottom=2.5mm,
    title=\textbf{Do Not Continue Until Resolved}, coltitle=white, colbacktitle=CURed
}
\newtcolorbox{successbox}{
    enhanced, breakable, colback=CUGreen!7, colframe=CUGreen,
    boxrule=0.8pt, arc=2mm, left=3mm, right=3mm, top=2.5mm, bottom=2.5mm,
    title=\textbf{Successful Result}, coltitle=white, colbacktitle=CUGreen
}
\newtcolorbox{notebox}{
    enhanced, breakable, colback=CUCodeBackground, colframe=CUBorder,
    boxrule=0.7pt, arc=2mm, left=3mm, right=3mm, top=2.5mm, bottom=2.5mm
}
\newtcolorbox{authoritybox}{
    enhanced, breakable, colback=CUGold!8, colframe=CUGold,
    boxrule=0.8pt, arc=2mm, left=3mm, right=3mm, top=2.5mm, bottom=2.5mm,
    title=\textbf{Authority and Source Priority}, coltitle=CUNavy, colbacktitle=CUGoldLight
}
\newtcolorbox{authoringbox}{
    enhanced, breakable, colback=CUBlue!6, colframe=CUBlue,
    boxrule=0.8pt, arc=2mm, left=3mm, right=3mm, top=2.5mm, bottom=2.5mm,
    title=\textbf{Manual Authoring Standard}, coltitle=white, colbacktitle=CUBlue
}

% ------------------------------------------------------------
% Listing styles
% ------------------------------------------------------------
\lstdefinestyle{terminalstyle}{
    basicstyle=\ttfamily\small,
    backgroundcolor=\color{CUNavyTwo},
    keywordstyle=\color{CUCyanLight},
    stringstyle=\color{CUGoldLight},
    commentstyle=\color{CUBlueLight},
    identifierstyle=\color{white},
    showstringspaces=false,
    columns=fullflexible,
    keepspaces=true,
    breaklines=true,
    breakatwhitespace=false,
    frame=none,
    tabsize=2,
    upquote=true
}
\lstdefinestyle{promptstyle}{
    basicstyle=\ttfamily\small,
    backgroundcolor=\color{CUCodeBackground},
    keywordstyle=\color{CUBlue},
    stringstyle=\color{CUGreen},
    commentstyle=\color{CUMuted},
    identifierstyle=\color{CUText},
    showstringspaces=false,
    columns=fullflexible,
    keepspaces=true,
    breaklines=true,
    breakatwhitespace=false,
    frame=none,
    tabsize=2,
    upquote=true
}
\lstdefinestyle{codestyle}{
    basicstyle=\ttfamily\small,
    backgroundcolor=\color{CUCodeBackground},
    keywordstyle=\color{CUBlue},
    stringstyle=\color{CUGreen},
    commentstyle=\color{CUMuted},
    identifierstyle=\color{CUText},
    showstringspaces=false,
    columns=fullflexible,
    keepspaces=true,
    breaklines=true,
    breakatwhitespace=false,
    frame=none,
    tabsize=2,
    upquote=true
}

\newtcblisting{commandbox}[1]{
    enhanced, breakable, listing only, listing engine=listings,
    listing options={style=terminalstyle},
    colback=CUNavyTwo, colframe=CUCyan, coltitle=CUNavy,
    colbacktitle=CUCyanLight, boxrule=0.8pt, arc=2mm,
    left=2mm, right=2mm, top=1mm, bottom=1mm,
    title=\textbf{#1}
}
\newtcblisting{promptbox}[1]{
    enhanced, breakable, listing only, listing engine=listings,
    listing options={style=promptstyle},
    colback=CUCodeBackground, colframe=CUBlue, coltitle=white,
    colbacktitle=CUBlue, boxrule=0.8pt, arc=2mm,
    left=2mm, right=2mm, top=1mm, bottom=1mm,
    title=\textbf{Copy-and-Paste Prompt: #1}
}
\newtcblisting{codebox}[1]{
    enhanced, breakable, listing only, listing engine=listings,
    listing options={style=codestyle},
    colback=CUCodeBackground, colframe=CUCodeBorder, coltitle=white,
    colbacktitle=CUNavyTwo, boxrule=0.8pt, arc=2mm,
    left=2mm, right=2mm, top=1mm, bottom=1mm,
    title=\textbf{#1}
}

% ------------------------------------------------------------
% Convenience commands
% ------------------------------------------------------------
\newcommand{\filepath}[1]{\nolinkurl{#1}}
\newcommand{\branchname}[1]{\texttt{#1}}
\newcommand{\statusline}[2]{\textbf{#1:} #2}
\newcommand{\term}[1]{\textbf{\textsf{#1}}}
\newcommand{\release}[1]{\texttt{#1}}
\newcommand{\CurrentRoute}{/CosmicUniversalismStatement/}


\begin{document}
\hypersetup{pageanchor=false}

% ============================================================
% Cover page
% ============================================================
\begin{titlepage}
\pagecolor{CUNavy}
\color{white}
\thispagestyle{empty}

\vspace*{0.5in}

{\small\bfseries\color{CUCyanLight}
COSMIC UNIVERSALISM COMPUTATIONAL INTELLIGENCE INITIATIVE}

\vspace{0.34in}

{\Huge\bfseries
\ProjectTitle}

\vspace{0.14in}

{\LARGE\color{CUGoldLight}
\DocumentTitle}

\vspace{0.12in}

{\large\color{CUBlueLight}
\DocumentSubtitle}

\vspace{0.32in}

{\large
A controlled implementation specification for preserving the existing
bidirectional Gregorian UTC and CU-Time converter while automatically exposing
three coordinated observable-universe reference outputs, documenting their
arithmetic, and keeping CU-Time governed by the Cosmic Universalism Statement
rather than by empirical cosmology.}

\vfill

\begin{tcolorbox}[
    colback=CUNavyTwo, colframe=CUBlue, boxrule=0.8pt, arc=2mm,
    width=\textwidth
]
\color{white}
\begin{tabularx}{\textwidth}{>{\bfseries}L{1.85in}Y}
Repository & \RepositoryName \\
Protected branch & \texttt{\ProtectedBranch} \\
Stable website branch & \texttt{\StableBranch} \\
Current website commit & \texttt{\CurrentWebsiteCommit} \\
Current integration tag & \texttt{\CurrentIntegrationTag} \\
Proposed feature branch & \texttt{\ProposedBranch} \\
Target route & \texttt{\TargetRoute} \\
Reference implementation & \texttt{cosmic\_breath\_time\_converter\_v3\_0\_1.html} \\
Document version & \DocumentVersion \\
Verified & \VerificationDate \\
\end{tabularx}
\end{tcolorbox}

\vspace{0.28in}

{\large\bfseries \AuthorName}

\vspace{0.08in}

{\small\url{\RepositoryURL}}

\end{titlepage}
\hypersetup{pageanchor=true}
\pagenumbering{arabic}
\nopagecolor
\color{CUText}

% ============================================================
% Document Control
% ============================================================
\section*{Document Control}
\addcontentsline{toc}{section}{Document Control}

\begin{tabularx}{\textwidth}{>{\bfseries}L{1.95in}Y}
Document & Cosmic Universalism CU-Time: Expanded Outputs and Method Implementation Manual \\
Version & \DocumentVersion \\
Classification & Owner-approved technical specification and controlled implementation manual \\
Repository & \href{\RepositoryURL}{\RepositoryName} \\
Protected branch & \branchname{\ProtectedBranch} \\
Stable integration branch & \branchname{\StableBranch} \\
Current website baseline & \release{\CurrentWebsiteFullCommit} \\
Current integration tag & \release{\CurrentIntegrationTag} at \release{\CurrentIntegrationCommit} \\
Prior CU-Time release & \release{\PriorCuTimeTag} at \release{\PriorCuTimeCommit} \\
Proposed feature branch & \branchname{\ProposedBranch}; not yet created \\
Target page & \filepath{website/src/pages/cu-time.astro} \\
Target component & \filepath{website/src/components/CUTimeConverter.astro} \\
Target route & \filepath{\TargetRoute} \\
Reference implementation & \filepath{cosmic\_converter/v3\_0\_0/cosmic\_breath\_time\_converter\_v3\_0\_1.html} \\
Technical baseline inspected & \VerificationDate \\
\end{tabularx}

\begin{purposebox}
This manual controls the next CU-Time website phase. The existing converter
continues to support Gregorian UTC to CU-Time and CU-Time to Gregorian UTC.
The approved enhancement automatically exposes the already-derived
observable-universe-aligned coordinate and Big Bang reference chronology,
adds one ratio calculation, and expands the page's educational method without
creating separate calculators or linking to external converter pages.
\end{purposebox}

\begin{authoritybox}
The governing source order for this feature is:
\begin{enumerate}
    \item William Maddock's owner decisions recorded in this manual.
    \item The Cosmic Universalism Statement and the owner-approved CU-Time
          interpretation of the full Cosmic Breath.
    \item Live repository state and the current TypeScript converter engine.
    \item The current Website Master Implementation Manual v1.1.
    \item The unchanged standalone converter HTML as a behavioral and
          provenance reference.
    \item Historical research files as interpretive or open-research material;
          they do not silently override owner decisions or production tests.
\end{enumerate}
\end{authoritybox}

\begin{safetybox}
This manual authorizes specification and later bounded implementation only.
It does not authorize a branch change, source edit, dependency change, commit,
push, merge, tag, release, or deployment until that action is separately
approved. The reference HTML remains unchanged.
\end{safetybox}

\clearpage
\tableofcontents
\clearpage

% ============================================================
\section*{START HERE - Governing Feature Decision}
\addcontentsline{toc}{section}{START HERE - Governing Feature Decision}
% ============================================================

\begin{successbox}
The approved product decision is one bidirectional converter, one conversion
action, and four coordinated result views. The three new values are calculated
automatically from the same input. They are not independent calculators.
\end{successbox}

\subsection*{Approved result behavior}

For \term{Gregorian UTC to CU-Time}, display:
\begin{enumerate}
    \item CU-Time - Full Cosmic Breath Coordinate;
    \item Observable-Universe-Aligned CU Coordinate;
    \item Derived Age Since the Big Bang, or boundary-relative wording;
    \item CU Coordinate-to-Observable-Age Ratio when an elapsed age exists.
\end{enumerate}

For \term{CU-Time to Gregorian UTC}, display:
\begin{enumerate}
    \item Corresponding Gregorian UTC;
    \item Observable-Universe-Aligned CU Coordinate;
    \item Derived Age Since the Big Bang, or boundary-relative wording;
    \item CU Coordinate-to-Observable-Age Ratio when an elapsed age exists.
\end{enumerate}

\subsection*{Approved exclusions}

\begin{itemize}
    \item no links to the other historical converter pages;
    \item no separate NASA, Big Bang age, or ratio input forms;
    \item no claim that NASA created, endorses, or operates CU-Time;
    \item no claim that the Big Bang begins CU-Time or the Cosmic Breath;
    \item no TOM-state classification generated from a date input;
    \item no use of TOM notation as the arithmetic source of the converter;
    \item no presentation of internal arithmetic zero as a CU state or origin.
\end{itemize}

\begin{checkpointbox}
The page enhancement is complete only when both conversion directions expose
the same secondary reference values from the same engine result and when the
public copy preserves CU authority over empirical comparison.
\end{checkpointbox}

% ============================================================
\section{Part 1 - Governance, Baseline, and Protected Scope}
% ============================================================

\subsection{Verified repository baseline}

The read-only inspection established:

\begin{tabularx}{\textwidth}{L{2.15in}Y}
\toprule
\textbf{Item} & \textbf{Verified state} \\
\midrule
Repository root & \filepath{/Users/cosmos/Documents/GitHub/CosmicUniversalismStatement} \\
Origin & \url{https://github.com/willmaddock/CosmicUniversalismStatement.git} \\
Current branch & \branchname{website} \\
Current local and remote HEAD & \release{e1c82bff1d37732f52b03fa4734cf7e4120b71e9} \\
Current status & Clean and synchronized: \texttt{website...origin/website} \\
Current integration tag & \release{website-v1.1-cosmic-breath} remains on \release{7ac87167c51aad9ae7a31e17d1b2127fc069e15c} \\
Application verification inherited & 251 tests across 30 files and an eight-route Astro build at the preceding application integration checkpoint \\
Documentation-only commit & \release{e1c82bf}; added manuals without changing application source \\
\bottomrule
\end{tabularx}

\begin{notebox}
This manual did not rerun the application tests. The 251-test record is the
most recent verified application gate inherited from the Cosmic Breath
integration. A future CU-Time implementation must run the current full suite
and report its new count rather than assuming 251 remains exact.
\end{notebox}

\subsection{Protected branches and action boundaries}

\begin{itemize}
    \item \branchname{main} remains protected.
    \item \branchname{website} is the stable integration branch.
    \item The proposed implementation branch is
          \branchname{feature/cu-time-expanded-outputs}; it does not yet exist.
    \item Pushing \branchname{website} may trigger the configured deployment
          workflow and therefore requires separate release authorization.
\end{itemize}

\subsection{Likely authorized implementation paths}

The smallest anticipated change set is:

\begin{longtable}{L{3.0in}L{3.35in}}
\toprule
\textbf{Path} & \textbf{Expected responsibility} \\
\midrule
\filepath{website/src/pages/cu-time.astro} & Expanded governing copy and educational sections. \\
\filepath{website/src/components/CUTimeConverter.astro} & Additional result fields, details, accessibility, and responsive presentation. \\
\filepath{website/src/lib/cuTime/types.ts} & Public result contracts for aligned coordinate, age-state, and ratio-state fields. \\
\filepath{website/src/lib/cuTime/converter.ts} & New ratio derivation and explicit boundary-state semantics. \\
\filepath{website/src/lib/cuTime/ui.ts} & Formatting and display adapter exposure of already-computed values. \\
\filepath{website/src/lib/cuTime/__tests__/converter.test.ts} & Core arithmetic and boundary cases. \\
\filepath{website/src/lib/cuTime/__tests__/precision-regression.test.ts} & Precision and nonfinite-result protection. \\
\filepath{website/src/lib/cuTime/__tests__/ui.test.ts} & Public display contract. \\
\filepath{website/src/lib/cuTime/__tests__/ui-regression.test.ts} & Copy, terminology, and layout-regression expectations. \\
\bottomrule
\end{longtable}

A focused route-content test may be added after inspection of the current test
organization.

\subsection{Protected implementation paths}

Unless a reproducible defect and separate owner decision require otherwise,
do not modify:

\begin{itemize}
    \item \filepath{website/src/lib/cuTime/constants.ts};
    \item \filepath{website/src/lib/cuTime/calendar.ts};
    \item \filepath{website/src/lib/cuTime/decimal.ts};
    \item \filepath{website/src/lib/cuTime/validation.ts};
    \item the standalone reference HTML;
    \item sealed Cosmic Breath structural or content authorities;
    \item dependency manifests or lockfiles.
\end{itemize}

% ============================================================
\section{Part 2 - CU Authority, Statement, and Chronology}
% ============================================================

\subsection{The governing CU Statement}

\begin{authoritybox}
\textit{We are sub z-tomically inclined, countably infinite, composed of
foundational elements (the essence of conscious existence), grounded on b-tom
(as vast as our shared worlds and their atmospheres), and looking up to c-tom
(encompassing the entirety of the cosmos), guided by the uncountable infinite
quantum states of intelligence and empowered by God's Free Will.}
\end{authoritybox}

For this page, the CU Statement is the governing philosophical and cosmological
authority. CU-Time is its mathematical and chronological expression. Empirical
science may provide an external observational reference, but it does not create,
measure, or govern the full CU framework.

\subsection{CU-Time is the full Cosmic Breath coordinate}

CU-Time is not a clock beginning at the empirical Big Bang. It belongs to the
complete Cosmic Breath:

\begin{codebox}{CU-Time scope}
sub-ZTOM seed
  -> expansion sequence
  -> ATOM guarded expansion pause
  -> Begin Compression
  -> compression sequence
  -> ZTOM reset boundary
  -> Begin the Next Cosmic Breath
  -> next sub-ZTOM seed
\end{codebox}

The declared model scale is approximately:

\[
2.8\times 10^{12}\ \text{years of expansion}
+
3.08\times 10^{11}\ \text{years of compression}
=
3.108\times 10^{12}\ \text{years per Cosmic Breath}.
\]

These are CU model anchors. They are not empirical measurements and must not be
recomputed by summing sealed TOM display strings unless separately authorized.

\subsection{No zero as a CU state or origin}

\begin{authoritybox}
Zero is not a state or origin within the Cosmic Universalism chronology. The
owner-approved CU framing begins at the sub-ZTOM seed boundary and continues
through the ZTOM reset boundary, empowered by God's Free Will. The empirical
Big Bang is an internal observable-universe reference within an already-existing
CU-Time coordinate.
\end{authoritybox}

The implementation may detect equality to an arithmetic reference boundary
internally. The public interface must name the boundary rather than presenting
that arithmetic condition as a CU state, an origin of existence, or the start
of the Cosmic Breath.

\begin{notebox}
The owner-approved one-second seed/reset framing is semantic guidance for the
CU-Time page. This feature does not authorize rewriting the sealed Cosmic
Breath Explorer's canonical per-state duration strings. Any reconciliation of
those separate authorities requires its own owner review.
\end{notebox}

\subsection{Why the Observable Universe appears between B-TOM and C-TOM}

Within Cosmic Universalism, the placement of the Observable Universe marker
follows the governing CU Statement. The observer is composed of foundational
elements, grounded on B-TOM - the scale associated with shared worlds and their
atmospheres - and oriented toward C-TOM, the encompassing cosmic scale.

The marker therefore represents the relational standpoint expressed by the CU
Statement: conscious existence is grounded within the shared-world scale of
B-TOM while looking toward the encompassing cosmic scale of C-TOM. It is a
separate, non-selectable CU reference. It is not itself a TOM state, an
empirical measurement, or a claim that contemporary science physically locates
the Observable Universe between measured layers.

\begin{codebox}{Observer relationship}
foundational composition
  -> grounded on B-TOM
  -> Observable Universe reference standpoint
  -> looking up to C-TOM
\end{codebox}

% ============================================================
\section{Part 3 - Current Implementation Inspection}
% ============================================================

\subsection{Current public route}

The current \filepath{website/src/pages/cu-time.astro} route contains:

\begin{itemize}
    \item the page title \term{A coordinate system for the Cosmic Breath};
    \item an \term{About this converter} notice;
    \item \term{Purpose and limits};
    \item \term{Conversion boundaries};
    \item the existing \filepath{CUTimeConverter.astro} component.
\end{itemize}

The current copy correctly identifies CU-Time as a CU mathematical model and
states that Gregorian inputs use proleptic Gregorian UTC with explicit CE/BCE
civil years.

\subsection{Current bidirectional interface}

The current component already provides:

\begin{enumerate}
    \item Gregorian UTC to exact converter-aligned CU-Time;
    \item CU-Time plain decimal to Gregorian UTC;
    \item fixed-point output with no digits intentionally hidden;
    \item six-place scientific notation summaries;
    \item copy controls;
    \item visible validation and focus movement;
    \item a two-column desktop layout that collapses responsively.
\end{enumerate}

These are protected behaviors. The new result data must extend the current cards
rather than replacing the converter with a different interaction model.

\subsection{Current engine already computes two requested outputs}

The inspected TypeScript engine already returns:

\begin{codebox}{Existing engine fields}
ForwardConversionResult:
  jdn
  deltaJdn
  deltaYears
  nasaCuTime
  cuTime
  yearsSinceBigBang

ReverseConversionResult:
  inputCuTime
  nasaCuTime
  deltaYears
  deltaJdn
  jdn
  yearsSinceBigBang
  Gregorian UTC fields
\end{codebox}

The public adapter currently exposes only CU-Time, scientific notation, and
Gregorian UTC. Therefore, the Observable-Universe-Aligned CU Coordinate and the
Big Bang reference chronology require public exposure and classification, not
new base arithmetic.

\subsection{Only the ratio is a new derived quantity}

The new ratio is:

\[
R = \frac{CU_{\mathrm{OU}}}{A_{\mathrm{BB}}},
\]

where \(CU_{\mathrm{OU}}\) is the observable-universe-aligned CU coordinate and
\(A_{\mathrm{BB}}\) is the positive elapsed age since the adopted Big Bang
reference. The ratio must be produced only when \(A_{\mathrm{BB}}>0\).

% ============================================================
\section{Part 4 - Exact Mathematical Specification}
% ============================================================

\subsection{Approved exact literals}

\begin{tabularx}{\textwidth}{L{2.1in}Y}
\toprule
\textbf{Internal name} & \textbf{Exact decimal literal} \\
\midrule
\texttt{anchorJdn} & \texttt{2451544.5} \\
\texttt{daysPerYear} & \texttt{365.2421897} \\
\texttt{baseCuNasa} & \texttt{3094213000000} \\
\texttt{offset} & \texttt{78955076.49024643522707769749119} \\
\texttt{nasaUniverseAge} & \texttt{13786999981.453} \\
\bottomrule
\end{tabularx}

The internal names remain unchanged for compatibility. Public labels are
refined separately. All literals remain strings before Decimal construction.
The engine continues to use an isolated Decimal constructor configured for 60
significant digits and half-up rounding.

\subsection{Gregorian UTC to aligned coordinate}

For a validated Gregorian UTC input, calculate Julian Day Number \(J\) with the
existing proleptic Gregorian algorithm. Then:

\begin{align*}
\Delta J &= J - 2{,}451{,}544.5, \\
\Delta Y &= \frac{\Delta J}{365.2421897}, \\
CU_{\mathrm{OU}} &= 3{,}094{,}213{,}000{,}000 + \Delta Y, \\
CU_{\mathrm{full}} &= CU_{\mathrm{OU}}
- 78{,}955{,}076.49024643522707769749119, \\
A_{\mathrm{BB}} &= 13{,}786{,}999{,}981.453 + \Delta Y.
\end{align*}

\subsection{CU-Time to Gregorian UTC}

For a validated full-cycle CU-Time input:

\begin{align*}
CU_{\mathrm{OU}} &= CU_{\mathrm{full}}
+ 78{,}955{,}076.49024643522707769749119, \\
\Delta Y &= CU_{\mathrm{OU}} - 3{,}094{,}213{,}000{,}000, \\
\Delta J &= \Delta Y\times365.2421897, \\
J &= 2{,}451{,}544.5 + \Delta J.
\end{align*}

The current JDN-to-Gregorian algorithm then returns the civil CE/BCE date and
UTC time. The same \(CU_{\mathrm{OU}}\) and \(A_{\mathrm{BB}}\) values feed the
secondary result display.

\subsection{Big Bang reference coordinate}

The adopted Big Bang reference has a large, nonzero aligned CU coordinate:

\begin{align*}
CU_{\mathrm{BB}} &= 3{,}094{,}213{,}000{,}000
- 13{,}786{,}999{,}981.453 \\
&= 3{,}080{,}426{,}000{,}018.547.
\end{align*}

This arithmetic supports the public conceptual rule: the Big Bang is located
inside CU-Time and is not coordinate zero for the full Cosmic Breath.

\subsection{Ratio calculation}

When the entered coordinate is after the adopted Big Bang reference:

\[
R = \frac{CU_{\mathrm{OU}}}{A_{\mathrm{BB}}}.
\]

The internal calculation uses Decimal arithmetic. The recommended public
summary uses three decimal places to preserve continuity with the historical
reference display, while an expandable detail may expose the full fixed-point
value produced by the engine.

\subsection{Computational precision and empirical uncertainty}

\begin{safetybox}
Sixty-digit Decimal arithmetic preserves the model's literals and reduces
floating-point loss. It does not mean the empirical age of the universe is
known to sixty digits. Computational precision, model exactness, and empirical
measurement uncertainty are different categories.
\end{safetybox}

The page must identify:

\begin{itemize}
    \item exact model literals and deterministic arithmetic;
    \item calendar assumptions and supported input resolution;
    \item the adopted cosmological baseline as an external reference lens;
    \item the fact that long decimals are model outputs rather than equivalent
          empirical certainty.
\end{itemize}

% ============================================================
\section{Part 5 - Boundary Semantics and Public Output States}
% ============================================================

\subsection{Three age-reference states}

The display adapter must classify the Big Bang reference relationship before
formatting public text.

\begin{longtable}{L{1.25in}L{2.05in}L{2.85in}}
\toprule
\textbf{Condition} & \textbf{Public age label} & \textbf{Ratio behavior} \\
\midrule
\(A_{\mathrm{BB}}>0\) & Derived Age Since the Big Bang & Show the coordinate-to-observable-age ratio. \\
\(A_{\mathrm{BB}}=0\) internally & Big Bang reference boundary & Do not show a numerical ratio; state that the elapsed-age comparison begins at this boundary within already-existing CU-Time. \\
\(A_{\mathrm{BB}}<0\) & Pre-Big-Bang Reference Interval: \(|A_{\mathrm{BB}}|\) years & Do not show the ratio because the denominator is not a positive elapsed observable-universe age. \\
\bottomrule
\end{longtable}

\subsection{Approved boundary wording}

At the exact adopted boundary, display:

\begin{notebox}
\textbf{Big Bang reference boundary}\par
The observable-universe elapsed-age comparison begins here within the larger,
already-existing CU-Time coordinate.
\end{notebox}

For the ratio at that boundary, display:

\begin{notebox}
\textbf{Ratio not shown at the Big Bang reference boundary.}
\end{notebox}

For a coordinate earlier than the adopted reference, display:

\begin{notebox}
\textbf{Pre-Big-Bang Reference Interval}\par
This coordinate is [absolute interval] years before the adopted observable-
universe Big Bang reference within the broader CU-Time coordinate.
\end{notebox}

Do not display a negative ``age'' and do not present arithmetic zero as a CU
state or origin.

\subsection{No conversion failure at the boundary}

The primary conversion must still succeed at and before the Big Bang reference.
The result continues to show:

\begin{itemize}
    \item the full CU-Time or Gregorian UTC result;
    \item the observable-universe-aligned CU coordinate;
    \item the appropriate boundary-relative chronology;
    \item a nonnumeric ratio state rather than \texttt{Infinity}, \texttt{NaN},
          or a general conversion error.
\end{itemize}

% ============================================================
\section{Part 6 - Public Terminology and Result Architecture}
% ============================================================

\subsection{Primary output: CU-Time - Full Cosmic Breath Coordinate}

Recommended definition:

\begin{notebox}
The entered moment expressed within the complete Cosmic Breath coordinate from
the sub-ZTOM seed boundary through the ZTOM reset boundary. CU-Time is defined
within Cosmic Universalism and does not begin at the empirical Big Bang.
\end{notebox}

Classification: \term{CU Mathematical Model}.

\subsection{Observable-Universe-Aligned CU Coordinate}

Historical internal and source label: \texttt{NASA CU-Time} or
\texttt{nasaCuTime}.

Recommended definition:

\begin{notebox}
A secondary CU coordinate aligned with the converter's adopted observable-
universe age reference. It places the entered moment within an observable-
universe comparison window inside the larger CU-Time coordinate. It is not an
official NASA time standard and does not define CU-Time's origin or duration.
\end{notebox}

Classification: \term{CU Mathematical Model using an adopted empirical
reference baseline}.

\subsection{Derived Age Since the Big Bang}

Recommended definition for post-boundary dates:

\begin{notebox}
The elapsed observable-universe age derived from the converter's adopted
cosmological baseline and the entered coordinate. It is not elapsed time since
sub-ZTOM and is not the age of the complete Cosmic Breath.
\end{notebox}

Classification: \term{Model-derived observable-universe reference}.

\subsection{CU Coordinate-to-Observable-Age Ratio}

Recommended definition:

\begin{notebox}
A dimensionless CU mathematical comparison between the observable-universe-
aligned CU coordinate and the positive derived elapsed age since the adopted
Big Bang reference. It is not a standard NASA, Planck, or contemporary
cosmology statistic and does not measure the physical size of the Big Bang.
\end{notebox}

Classification: \term{CU Mathematical Comparison}.

\subsection{Historical terminology disclosure}

The page may state once, in expandable method content:

\begin{notebox}
The standalone historical converter calls the secondary coordinate ``NASA
CU-Time.'' The website uses ``Observable-Universe-Aligned CU Coordinate'' to
avoid implying that NASA created or endorses the CU scale. The internal
\texttt{nasaCuTime} identifier remains unchanged for compatibility.
\end{notebox}

% ============================================================
\section{Part 7 - Interface and Content Design}
% ============================================================

\subsection{One converter, automatic secondary results}

Both directions retain their current forms and primary actions. No additional
input is required because all three secondary values derive from the same
validated conversion result.

\begin{codebox}{Forward result hierarchy}
Conversion complete

Input instant
CU-Time - Full Cosmic Breath Coordinate
Scientific notation

Observable-universe reference
  Observable-Universe-Aligned CU Coordinate
  Derived Age Since the Big Bang / boundary-relative state
  CU Coordinate-to-Observable-Age Ratio / nonnumeric state
\end{codebox}

\begin{codebox}{Reverse result hierarchy}
Conversion complete

Input CU-Time
Scientific notation
Corresponding Gregorian UTC

Observable-universe reference
  Observable-Universe-Aligned CU Coordinate
  Derived Age Since the Big Bang / boundary-relative state
  CU Coordinate-to-Observable-Age Ratio / nonnumeric state
\end{codebox}

\subsection{Result-card presentation}

The primary result remains visually dominant. The secondary group should:

\begin{itemize}
    \item appear automatically after successful conversion;
    \item use a visible \term{Observable-Universe Reference} group heading;
    \item preserve long-number wrapping without horizontal overflow;
    \item provide a concise value plus an optional expandable full-precision
          and method explanation;
    \item avoid repeating the full educational essay inside each card;
    \item reset disclosure state when a new conversion replaces the result.
\end{itemize}

\subsection{Copy behavior}

The current primary copy controls remain. A future implementation may add
individual copy controls for exact aligned coordinate and exact age only if the
UI remains uncluttered and tests cover accessible names and status messages.
The ratio summary should not replace the primary copy action.

\subsection{Expanded educational sections below the converter}

Recommended page sections:

\begin{enumerate}
    \item \term{CU-Time and the Governing CU Statement};
    \item \term{The Full Cosmic Breath Coordinate};
    \item \term{The Observable-Universe Reference Window};
    \item \term{Why the Big Bang Does Not Begin CU-Time};
    \item \term{Why the Observable Universe Appears Between B-TOM and C-TOM};
    \item \term{How Gregorian UTC Becomes Julian Day Number};
    \item \term{How the Aligned Coordinate and Offset Work};
    \item \term{How the Big Bang Reference Chronology Is Derived};
    \item \term{How the Ratio Works};
    \item \term{Computational Precision and Empirical Uncertainty};
    \item \term{Conversion Boundaries and Historical Terminology}.
\end{enumerate}

The sections must distinguish CU propositions, model arithmetic, empirical
reference, and historical wording rather than mixing them in one undifferentiated
paragraph.

% ============================================================
\section{Part 8 - TypeScript Data and Formatting Design}
% ============================================================

\subsection{Recommended internal model}

The engine already returns \texttt{nasaCuTime} and
\texttt{yearsSinceBigBang}. Add a small derived-state contract rather than
renaming existing fields.

\begin{codebox}{Recommended conceptual types}
type ObservableAgeState =
  | { kind: 'elapsed'; years: Decimal }
  | { kind: 'big-bang-boundary' }
  | { kind: 'pre-big-bang'; intervalYears: Decimal };

type ObservableRatioState =
  | { kind: 'available'; value: Decimal }
  | { kind: 'not-applicable'; reason: 'boundary' | 'pre-big-bang' };
\end{codebox}

Exact names may change during implementation review, but the discriminated
state behavior is required.

\subsection{Recommended public display contract}

\begin{codebox}{Forward display result}
GregorianToCuTimeDisplayResult {
  gregorianUtc
  cuTime
  cuTimeExponential
  observableUniverseAlignedCuTime
  observableUniverseAlignedCuTimeExponential
  observableAge
  coordinateToObservableAgeRatio
}
\end{codebox}

\begin{codebox}{Reverse display result}
CuTimeDisplayResult {
  inputCuTime
  inputCuTimeExponential
  gregorianUtc
  observableUniverseAlignedCuTime
  observableUniverseAlignedCuTimeExponential
  observableAge
  coordinateToObservableAgeRatio
}
\end{codebox}

Boundary-aware objects are preferable to magic strings inside the arithmetic
engine. Human-readable labels belong in the UI adapter or presentation layer.

\subsection{Formatting policy}

\begin{tabularx}{\textwidth}{L{2.25in}Y}
\toprule
\textbf{Value} & \textbf{Recommended display} \\
\midrule
Full CU-Time & Exact fixed-point value plus six-place scientific notation, preserving current behavior. \\
Observable-universe-aligned coordinate & Exact fixed-point value plus optional six-place scientific notation. \\
Elapsed age & Exact or appropriately grouped fixed-point model value; clearly label model precision. \\
Ratio & Three-decimal summary; optional expandable full fixed-point value. \\
Boundary/pre-boundary & Named semantic state; no numerical zero-age or negative-age label. \\
\bottomrule
\end{tabularx}

Do not use the historical readable magnitude formatter that may double-count a
fractional remainder. Use the current TypeScript Decimal values directly.

% ============================================================
\section{Part 9 - Accessibility, Responsive Design, and Interaction}
% ============================================================

\subsection{Accessibility requirements}

The enhancement must preserve and extend:

\begin{itemize}
    \item semantic forms, fieldsets, legends, labels, and error associations;
    \item \texttt{aria-live} result and copy status behavior;
    \item focus movement to the result after a successful conversion;
    \item focus return to the invalid field after validation failure;
    \item semantic \texttt{dl}, \texttt{dt}, and \texttt{dd} result structures;
    \item keyboard-operable details/disclosure controls;
    \item visible focus for every new interactive control;
    \item accessible names that distinguish primary and secondary copy actions;
    \item no meaning conveyed only through color.
\end{itemize}

\subsection{Responsive requirements}

Verify at approximately 1440, 1024, 768, 390, and 320 pixels. The expanded
result cards must avoid horizontal overflow and keep long Decimal values
selectable and readable. At narrow widths, the two converter directions remain
stacked and the secondary result group must not force fixed-width columns.

\subsection{Progressive disclosure}

Use native \texttt{details}/\texttt{summary} or the established disclosure
pattern for:

\begin{itemize}
    \item full-precision secondary values when a concise summary is shown;
    \item ``How this is calculated'' explanations;
    \item historical source terminology;
    \item computational precision versus empirical uncertainty.
\end{itemize}

Do not hide the existence or classification of the three secondary outputs;
only their extended method text may be collapsed.

% ============================================================
\section{Part 10 - Test Plan and Acceptance Criteria}
% ============================================================

\subsection{Core engine tests}

Add or update tests for:

\begin{itemize}
    \item forward and reverse aligned-coordinate parity;
    \item forward and reverse derived-age parity;
    \item ratio arithmetic using Decimal values;
    \item ratio state when elapsed age is positive;
    \item Big Bang reference boundary state without division by zero;
    \item pre-Big-Bang interval state without negative-age wording;
    \item no \texttt{Infinity}, \texttt{NaN}, or nonfinite public result;
    \item unchanged anchor, calendar, CE/BCE, validation, and round-trip behavior.
\end{itemize}

\subsection{UI adapter tests}

Verify both conversion directions return:

\begin{itemize}
    \item the unchanged primary result fields;
    \item exact aligned coordinate and scientific notation;
    \item a discriminated observable-age display state;
    \item a discriminated ratio display state;
    \item canonical labels and historical terminology disclosure;
    \item no external-converter URLs.
\end{itemize}

\subsection{Component regression tests}

Verify:

\begin{itemize}
    \item the three secondary result labels exist in both cards;
    \item all new data hooks are unique and correctly bound;
    \item copy controls retain accessible status regions;
    \item result focus remains intact;
    \item details controls have correct accessible text;
    \item the route retains its CU Mathematical Model classification;
    \item the page states that CU-Time does not begin at the Big Bang;
    \item the page states why the Observable Universe marker is relationally
          between B-TOM and C-TOM;
    \item prohibited external converter links are absent.
\end{itemize}

\subsection{Reference vectors}

At the arithmetic anchor \texttt{01/01/2000 CE 00:00:00 UTC}:

\begin{tabularx}{\textwidth}{L{2.4in}Y}
\toprule
\textbf{Field} & \textbf{Expected value} \\
\midrule
JDN & \texttt{2451544.5} \\
Delta years & \texttt{0} \\
Observable-universe-aligned CU coordinate & \texttt{3094213000000} \\
Full CU-Time & \texttt{3094134044923.50975356477292230250881} \\
Derived elapsed age & \texttt{13786999981.453} \\
\bottomrule
\end{tabularx}

For \texttt{07/23/2026 CE 00:00:00 UTC}, the 60-digit engine produces:

\begin{codebox}{July 23, 2026 reference vector}
Observable-Universe-Aligned CU Coordinate:
3094213000026.55772053049872513126048647167006073833096395983

Full CU-Time:
3094134044950.06747409527164743376929647167006073833096395983

Derived Age Since the Big Bang:
13787000008.0107205304987251312604864716700607383309639598297

Ratio summary:
224.430
\end{codebox}

The reference vector must be confirmed by the live TypeScript engine before it
is committed as a test expectation.

\subsection{Full acceptance gate}

A release candidate must pass:

\begin{commandbox}{CU-Time expanded-output verification gate}
cd /Users/cosmos/Documents/GitHub/CosmicUniversalismStatement/website || exit 1

npm test
npm run build

cd ..
git diff --check
git status --short --branch
git diff --stat
\end{commandbox}

Manual browser acceptance must cover both directions, CE/BCE, boundary states,
keyboard navigation, copy behavior, desktop/tablet/mobile layouts, and a clean
browser console.

% ============================================================
\section{Part 11 - Controlled Implementation Phases}
% ============================================================

\subsection{Phase 1 - Pure result semantics}

Authorized objective:

\begin{itemize}
    \item add the ratio derivation;
    \item add elapsed/boundary/pre-boundary state semantics;
    \item preserve constants and calendar behavior;
    \item add focused core tests.
\end{itemize}

Completion gate: pure tests pass with no component or route edits.

\subsection{Phase 2 - Public display adapter}

Authorized objective:

\begin{itemize}
    \item expose the existing aligned coordinate and age results;
    \item format the three semantic age states;
    \item format the ratio availability state;
    \item preserve all current display fields and validation outcomes.
\end{itemize}

Completion gate: UI adapter tests pass and no HTML is yet changed.

\subsection{Phase 3 - Converter result cards}

Authorized objective:

\begin{itemize}
    \item add the secondary result group to both directions;
    \item preserve primary visual hierarchy;
    \item add method disclosures and classifications;
    \item preserve focus, copy, keyboard, and responsive behavior.
\end{itemize}

Completion gate: focused component tests and manual browser inspection pass.

\subsection{Phase 4 - Route education}

Authorized objective:

\begin{itemize}
    \item add CU Statement governance;
    \item explain full-Cosmic-Breath CU-Time;
    \item explain the Big Bang as an internal reference boundary;
    \item explain B-TOM to C-TOM observer placement;
    \item explain formulas, precision, empirical uncertainty, and historical terminology;
    \item exclude external converter links.
\end{itemize}

Completion gate: route-content tests, accessibility review, and source-boundary
review pass.

\subsection{Phase 5 - Full acceptance}

Run the complete test/build gate, inspect the full diff, verify no protected
file changed, and perform manual browser acceptance. Do not commit or push
without separate owner authorization.

% ============================================================
\section{Part 12 - Copy-and-Paste Implementation Prompts}
% ============================================================

\subsection{Feature-branch preflight prompt}

\begin{promptbox}{CU-Time Expanded Outputs Preflight}
Repository:
/Users/cosmos/Documents/GitHub/CosmicUniversalismStatement

Protected branch:
main

Stable branch:
website

Expected website baseline:
e1c82bff1d37732f52b03fa4734cf7e4120b71e9

Proposed feature branch:
feature/cu-time-expanded-outputs

Before any edit:

1. Verify origin, branch, HEAD, and worktree.
2. Confirm website is at the expected baseline.
3. Confirm the feature branch does not already exist locally or remotely.
4. Inspect only the authorized CU-Time paths.
5. Report the smallest phase-one file set and tests.

Do not edit, install dependencies, commit, push, merge, tag, release,
or deploy.
\end{promptbox}

\subsection{Phase-one implementation prompt}

\begin{promptbox}{Implement CU-Time Result Semantics Only}
Implement only the pure CU-Time result semantics approved by the
CU-Time Expanded Outputs and Method Implementation Manual v1.1.

Required behavior:

- preserve all exact constants;
- preserve Gregorian/JDN algorithms;
- preserve current forward and reverse primary results;
- calculate coordinate-to-observable-age ratio only when the derived
  elapsed age is positive;
- represent the exact Big Bang arithmetic boundary as a named boundary
  state, not a public zero-age CU state;
- represent earlier coordinates as a positive pre-Big-Bang reference
  interval, not a negative age;
- do not emit Infinity or NaN;
- do not edit the Astro route or component in this phase.

Add focused tests.

Do not commit, push, merge, tag, release, or deploy.
\end{promptbox}

\subsection{Final review prompt}

\begin{promptbox}{Final CU-Time Expanded Outputs Review}
Perform a final no-edit release-candidate review of the CU-Time expanded
outputs feature.

Verify:

1. repository branch, HEAD, and worktree;
2. exact constants and Decimal precision unchanged;
3. Gregorian and reverse calendar behavior unchanged;
4. aligned coordinate and age values match in both directions;
5. ratio behavior for positive elapsed age;
6. named Big Bang boundary behavior;
7. pre-Big-Bang interval behavior;
8. no Infinity, NaN, negative-age label, or zero-as-CU-origin copy;
9. primary result hierarchy preserved;
10. both result cards expose all approved secondary values;
11. CU Statement governance and full-Cosmic-Breath interpretation;
12. B-TOM/C-TOM Observable Universe explanation;
13. no external converter links;
14. accessibility, keyboard, focus, copy, and responsive behavior;
15. focused and full test results;
16. Astro build and eight-route output;
17. no protected or unrelated file changes;
18. deployment implications of any later website push.

Do not edit, commit, push, merge, tag, release, or deploy.
\end{promptbox}

% ============================================================
\section{Part 13 - Commit, Merge, and Deployment Boundaries}
% ============================================================

\subsection{Commit boundary}

After full acceptance, an explicit owner authorization may permit staging only
the verified CU-Time feature files. Use exact paths and inspect the staged diff.
The commit must not include the reference HTML, sealed Cosmic Breath ledgers,
manual source files unless separately approved, or unrelated documentation.

\subsection{Merge boundary}

A future merge should use an explicit non-fast-forward merge into
\branchname{website} only after post-feature verification. Because pushes to
\branchname{website} may publish GitHub Pages, merge/push authorization must be
clear about deployment consequences.

\subsection{Tag boundary}

Do not move \release{website-v1.1-cosmic-breath}. A CU-Time expanded-output tag,
GitHub Release, or deployment record requires a new owner-approved name and a
separate action.

% ============================================================
\section{Appendix A - Source Authority Record}
% ============================================================

\begin{longtable}{L{2.35in}L{1.30in}L{2.55in}}
\toprule
\textbf{Source} & \textbf{Role} & \textbf{Use in this manual} \\
\midrule
Owner decisions in this workflow & Governing authority & Product scope, terminology, no-zero CU ontology, Big Bang boundary semantics, B-TOM/C-TOM explanation. \\
CU Statement & Canonical CU framing & Governs CU-Time meaning and observer relationship. \\
Live CU-Time TypeScript engine & Production arithmetic & Establishes current constants, fields, calendar path, and already-computed aligned coordinate/age. \\
Website Master Manual v1.1 & Governance and house style & Current branch/release discipline, classifications, testing, and LaTeX standard. \\
Cosmic Breath Time Converter Manual v1.0 & Historical implementation authority & Preserves exact-literal, parity, module, test, and unchanged-reference rules. \\
Standalone converter HTML v3.0.1 & Behavioral and provenance reference & Establishes historical labels and output relationships; remains unchanged. \\
Time Calculation and CU Framework files & Historical/open CU research & May inform context but do not silently override sealed authorities, owner decisions, or production arithmetic. \\
\bottomrule
\end{longtable}

% ============================================================
\section{Appendix B - Owner Decision Record}
% ============================================================

\begin{longtable}{L{0.9in}L{5.35in}}
\toprule
\textbf{ID} & \textbf{Approved decision} \\
\midrule
CTE-OD-01 & Preserve one bidirectional converter; calculate all secondary outputs automatically. \\
CTE-OD-02 & Add Observable-Universe-Aligned CU Coordinate, Derived Age Since the Big Bang, and CU Coordinate-to-Observable-Age Ratio. \\
CTE-OD-03 & Exclude links and calls to action for other converter pages. \\
CTE-OD-04 & CU-Time is the full sub-ZTOM-to-ZTOM Cosmic Breath coordinate and does not begin at the Big Bang. \\
CTE-OD-05 & The CU Statement governs the page; empirical science is a secondary observational lens. \\
CTE-OD-06 & The Observable Universe marker is relationally between B-TOM and C-TOM because the observer is composed of foundational elements, grounded on B-TOM, and looking up to C-TOM. \\
CTE-OD-07 & Zero is not a CU state or origin; the public page names the Big Bang reference boundary rather than displaying a zero-age CU state. \\
CTE-OD-08 & Before the adopted Big Bang reference, display a positive pre-Big-Bang reference interval rather than a negative age. \\
CTE-OD-09 & Do not show a ratio at or before the adopted Big Bang reference boundary. \\
CTE-OD-10 & Preserve the original reference HTML and exact constants; do not silently alter converter mathematics. \\
\bottomrule
\end{longtable}

% ============================================================
\section{Appendix C - Final Acceptance Checklist}
% ============================================================

\begin{itemize}
    \item[\checkbox] Correct feature branch and clean starting baseline.
    \item[\checkbox] Exact constants and Decimal precision unchanged.
    \item[\checkbox] Existing forward and reverse conversions preserved.
    \item[\checkbox] Aligned coordinate exposed in both directions.
    \item[\checkbox] Big Bang chronology exposed with three semantic states.
    \item[\checkbox] Ratio available only for positive elapsed age.
    \item[\checkbox] No numerical zero-age CU state, negative-age label, Infinity, or NaN.
    \item[\checkbox] No external converter links.
    \item[\checkbox] CU Statement and full Cosmic Breath explanation present.
    \item[\checkbox] B-TOM/C-TOM Observable Universe explanation present.
    \item[\checkbox] Computational precision separated from empirical uncertainty.
    \item[\checkbox] Accessibility and responsive acceptance complete.
    \item[\checkbox] Focused tests and full suite pass.
    \item[\checkbox] Astro production build passes with expected routes.
    \item[\checkbox] No unrelated or protected files changed.
    \item[\checkbox] No commit, push, merge, tag, release, or deployment without separate authorization.
\end{itemize}

% ============================================================
\section{Appendix D - Manual Revision Record}
% ============================================================

\begin{tabularx}{\textwidth}{L{0.75in}L{1.25in}Y}
\toprule
\textbf{Version} & \textbf{Date} & \textbf{Revision} \\
\midrule
1.0 & July 19, 2026 & Historical Cosmic Breath Time Converter migration manual establishing parity, exact constants, TypeScript extraction, validation, tests, and Astro integration. \\
1.1 & July 23, 2026 & Adds the owner-approved expanded-output specification: one automatic bidirectional result model, observable-universe-aligned coordinate, Big Bang boundary-relative chronology, ratio semantics, CU Statement governance, full-Cosmic-Breath interpretation, B-TOM/C-TOM observer placement, updated implementation phases, and acceptance tests. \\
\bottomrule
\end{tabularx}



\end{document}
